
The morphism \(\widetilde d_1\) is therefore quasi-finite and flat at
the points of \(F_1\) lying over \(z\). Since \(\widetilde d_1\) is
of finite type, SGA~1, IV, 6.10 and EGA~III, 4.4.10 imply that the set
\(\Phi\) of points of \(F_1\) at which \(\widetilde d_1\) is not both
flat and quasi-finite is closed in \(F_1\), and hence in \(X_1\)
because \(F_1\) is closed in \(X_1\).%
\footnote{\label{note:sga3-V-40}Editor's note: If \(f:X\to Y\) is
locally of finite type, the set of points \(x\in X\) that are isolated
in their fiber \(f^{-1}(f(x))\) is open in \(X\). In EGA~III, 4.4.10,
this is deduced from Zariski's Main Theorem when \(Y\) is locally
Noetherian. For arbitrary \(Y\), it follows from Chevalley's
semicontinuity theorem (EGA~IV\(_3\), 13.1.3 and 13.1.4).
Consequently, \(f\) is quasi-finite at \(x\) if and only if it is of
finite type at \(x\) and \(x\) is isolated in
\(f^{-1}(f(x))\). This will be used below; see \hyperref[note:sga3-V-42]{editor's note~(42)}.}
Since \(d_1\) is proper, \(\widetilde d_1(\Phi)\) is closed and, by
what precedes, does not contain \(z\). Put
\[
  W=\widetilde d_1(F_1)-\widetilde d_1(\Phi).
\]
The restriction of \(\widetilde d_1\) to
\(\widetilde d_1^{-1}(W)\) is of finite presentation (by the
Noetherian hypotheses), flat, proper, and quasi-finite.%
\footnote{\label{note:sga3-V-41}Editor's note: The original has been
expanded in what follows.}
It is therefore finite, locally free, and open, by EGA~III, 4.4.2 and
EGA~IV\(_2\), 2.1.12 and 2.4.6. Consequently,
\(\widetilde d_1(F_1)\) is a neighborhood of \(z\), and \(W\) is the
largest open subset of \(X_0\) contained in \(\widetilde d_1(F_1)\)
over which \(\widetilde d_1\) is both flat and quasi-finite.

We shall see in Lemma~\hyperref[lem:sga3-V-7.4]{7.4} that the inverse
images of \(\Phi\) under \(\widetilde d'_1\) and
\(\widetilde d'_0\) both identify with the set of points of \(F_2\)
at which \(\widetilde d'_2\) is not both flat and quasi-finite.
It follows that
\[
\begin{aligned}
  d_0^{-1}(W)
  &=
  \widetilde d'_2(F_2)
  -
  \widetilde d'_2\bigl((\widetilde d'_0)^{-1}(\Phi)\bigr),\\
  d_1^{-1}(W)
  &=
  \widetilde d'_2(F_2)
  -
  \widetilde d'_2\bigl((\widetilde d'_1)^{-1}(\Phi)\bigr)
\end{aligned}
\]
coincide. Thus \(W\) is saturated.

Set
\[
  W_1=\widetilde d_1^{-1}(W).
\]
The equality \(d_0^{-1}(W)=d_1^{-1}(W)\) implies
\[
  (\widetilde d'_2)^{-1}\bigl(d_0^{-1}(W)\bigr)
  =
  (\widetilde d'_2)^{-1}\bigl(d_1^{-1}(W)\bigr),
\]
that is,
\[
  (\widetilde d'_0)^{-1}(W_1)
  =
  (\widetilde d'_1)^{-1}(W_1).
\]
The morphism \(\widetilde d_0\) is faithfully flat and quasi-compact,
because \(d_0\), like \(d_1\), is surjective, proper, and flat.
Moreover, the square
\SGAThreeDiagram[0.30\textwidth]
  {figures/exp5/p277_lemma73_descent_square.png}
  {The Cartesian square used to descend \(W_1\).}
  {fig:sga3-V-7.3-descent-square}
  {Expos\'e V, component-local p.~27 / combined-reader p.~289 / re-edited header p.~277}
is Cartesian. It follows that
\[
  W_1=\widetilde d_0^{-1}(U)
\]
for an open subset \(U\) of \(F\) (SGA~1, VIII, \SGARefLink{sga3:viii:corollary:4:4}{4.4}). This open
subset \(U\) is a quasi-section of the groupoid with base \(W\)
induced by \(X_\ast\). We may therefore take \(U_z=W\).

It remains to state \hyperref[lem:sga3-V-7.4]{the following lemma}, whose proof is classical.

\medskip
\phantomsection\label{lem:sga3-V-7.4}%
\noindent\textbf{Lemma 7.4.}---%
\emph{Consider a Cartesian square of schemes}
\SGAThreeDiagram[0.30\textwidth]
  {figures/exp5/p277_lemma74_cartesian_square.png}
  {The Cartesian square in \hyperref[lem:sga3-V-7.4]{Lemma~7.4}.}
  {fig:sga3-V-7.4-cartesian-square}
  {Expos\'e V, component-local p.~27 / combined-reader p.~289 / re-edited header p.~277}
